\subsection{可积集}

定义 2. --- 局部紧空间 X 的子集 A 称为关于 X 上的测度 \(\mu\) 可积的（或称它是 \(\mu\) -可积的），如果 A 的特征函数 \(\varphi_{A}\) 可积。有限数 \(\mu(A) = \int \varphi_{A} d\mu\) 称为 A 的测度。

对每个可积集 A，\(|\mu|(A) = |\mu|^{*}(A)\) （No.~2, 命题 1）；为使一个集是可忽略的，必须且只须它关于 \(|\mu|\) 的测度为零。

命题 6. --- 可积集的有限族 \((\mathrm{A}_{i})_{1\leqslant i\leqslant n}\) 的并集是可积的，且

\[
| \mu | \left(\bigcup_ {i = 1} ^ {n} \mathrm{A} _ {i}\right) \leqslant \sum_ {i = 1} ^ {n} | \mu | (\mathrm{A} _ {i}).\tag{13}
\]

此外，若诸 \(\mathrm{A}_i\) 两两不相交，则

\[
\mu \left(\bigcup_ {i = 1} ^ {n} \mathrm{A} _ {i}\right) = \sum_ {i = 1} ^ {n} \mu (\mathrm{A} _ {i}).\tag{14}
\]

因为，若 \(\mathbf{A} = \bigcup_{i=1}^{n} \mathbf{A}_i\) ，则 \(\varphi_{\mathbf{A}} = \sup \varphi_{\mathbf{A}_i}\) ，因此（§3, No.~5, 命题 12 的推论）若诸 \(\mathbf{A}_i\) 可积，则 \(\mathbf{A}\) 也可积；关系式 (13) 是外测度的类似关系式（§1, No.~4, 命题 18）的特例，只需注意到关系式 \(|\mu|(\mathrm{A}) = |\mu|^*(\mathrm{A})\) ；最后，若诸 \(\mathbf{A}_i\) 两两不相交，则 \(\varphi_{\mathrm{A}} = \sum_{i=1}^{n} \varphi_{\mathrm{A}_i}\) ，由此得 (14)。

命题 7. --- \(1^{\circ}\) 若 A 与 B 是两个可积集，使 \(\mathbf{B} \subset \mathbf{A}\) ，则集 \(\mathbf{C} = \mathbf{A} - \mathbf{B}\) 是可积的，且

\[
\mu (\mathrm{C}) = \mu (\mathrm{A}) - \mu (\mathrm{B}).\tag{15}
\]

\(2^{\circ}\) 可积集的可数族的交是可积的。第一部分由 \(\varphi_{\mathrm{C}} = \varphi_{\mathrm{A}} - \varphi_{\mathrm{B}}\) 这一事实得出。另一方面，若 \((\mathbf{A}_n)\) 是可积集的一个序列而 A 是它们的交，则 \(\varphi_{\mathrm{A}} = \inf_{n} \varphi_{\mathrm{A}_n}\) ，因此 A 是可积的（No.~3, 命题 4）。

推论. --- 若 \((\mathrm{A}_{n})\) 是可积集的一个递减序列，则 \(\mu(\bigcap_{n}\mathrm{A}_{n})=\lim_{n\to\infty}\mu(\mathrm{A}_{n})\) 。

因为，若 \(\mathrm{A} = \bigcap_{n}\mathrm{A}_{n}\) ，则 \(\varphi_{\mathrm{A}}\) 是递减序列 \((\varphi_{\mathrm{A}_n})\) 的下包络（No.~3, 命题 4）。

命题 8. --- 设 \((\mathbf{A}_{n})\) 是可积集的一个递增序列；为使并集 \(A = \bigcup_{n} A_{n}\) 可积，必须且只须 \(\sup_{n} |\mu|(\mathbf{A}_{n}) < +\infty\) ；此时

\[
\mu (\mathrm{A}) = \lim _ {n \rightarrow \infty} \mu (\mathrm{A} _ {n}).\tag{16}
\]

因为，诸 \(\varphi_{\mathrm{A}_n}\) 构成可积函数的一个递增序列，且 \(\varphi_{\mathrm{A}} = \sup_n \varphi_{\mathrm{A}_n}\) ；于是命题由 No.~3 的命题 4 得出。

推论. --- 设 \((\mathrm{A}_{n})\) 是可积集的一个序列，使 \(\sum_{n=1}^{\infty} |\mu|(\mathrm{A}_{n}) < +\infty\) ；则并集 \(A = \bigcup_{n} A_{n}\) 是可积的，且

\[
| \mu | \left(\bigcup_ {n} \mathrm{A} _ {n}\right) \leqslant \sum_ {n = 1} ^ {\infty} | \mu | (\mathrm{A} _ {n}).\tag{17}
\]

因为，\(\varphi_{\mathrm{A}} = \sup_n\varphi_{\mathrm{A}_n}\) ，且

\[
| \mu | ^ {*} (\mathrm{A}) \leqslant \sum_ {n = 1} ^ {\infty} | \mu | ^ {*} (\mathrm{A} _ {n}) = \sum_ {n = 1} ^ {\infty} | \mu | (\mathrm{A} _ {n}) <   + \infty
\]

（§1, No.~4, 命题 18）；因此 A 是可积的（§3, No.~6, 定理 5 的推论 2），且由于 \(|\mu|(\mathrm{A}) = |\mu|^*(\mathrm{A})\) ，我们确有 (17)。

命题 9. --- 设 \((\mathbf{A}_{n})\) 是两两不相交的可积集的一个序列，使 \(\sum_{n=1}^{\infty}|\mu|(\mathbf{A}_{n})<+\infty\) ；则

\[
\mu \left(\bigcup_ {n} \mathrm{A} _ {n}\right) = \sum_ {n = 1} ^ {\infty} \mu (\mathrm{A} _ {n}).\tag{18}
\]

因为，若 \(\mathbf{A} = \bigcup_{n}\mathbf{A}_{n}\) ，则 \(\varphi_{\mathrm{A}} = \sum_{n = 1}^{\infty}\varphi_{\mathrm{A}_n}\) ，于是命题由 (17) 与 No.~3 定理 2 的推论 2 得出。

关系式 (18) 也可表述为：测度 \(\mu\) 在 X 的可积子集所成的集合上是完全可加的。

